%%=============================================================================
%% Proof of Concept
%%=============================================================================

\chapter{Proof of Concept}%
\label{ch:proofofconcept}

De proof of concept bevat een gedetailleerde bespreking van de implementatie van het model, inclusief de vergelijking tussen de verschillende versies van de code, van high-level tot low-level en \glsxtrshort{tflite}. Het doel is om de effectiviteit en efficiëntie van \gls{robbert} op verschillende platformen te evalueren en om de capaciteiten van een niet-gefinetunde \gls{robbert} te concluderen. In elke integratie die volgt, zijn de volgende modellen getest of gebruikt:

\begin{itemize}
    \item DTAI-KULeuven/robbertje-1-gb-non-shuffled5
    \item pdelobelle/robbert-v2-dutch-base
    \item DTAI-KULeuven/robbert-2022-dutch-base
    \item DTAI-KULeuven/robbert-2023-dutch-base
    \item DTAI-KULeuven/robbert-2023-dutch-large
\end{itemize}

\section{Python Implementatie}

Alle code die in deze sectie wordt besproken, is beschikbaar in de volgende GitHub-repository: \url{https://github.com/Chmpy/RobBERT_local_example}. Deze repository bevat de volledige implementatie van de scripts en configuraties die nodig zijn om de \glsxtrshort{tflite}-modellen te verkrijgen en gebruiken

\subsection{Vereisten en Setup}

In deze sectie worden de vereiste bibliotheken en hun versies besproken voor de implementatie van de proof of concept in Python. Daarnaast beschrijven we het utility-bestand dat zorgt voor consistent gedrag tussen de scripts.

Voor het uitvoeren van de scripts en het installeren van de benodigde bibliotheken wordt Python versie \texttt{3.11.5} gebruikt. Het is van cruciaal belang om deze versie te hanteren om ervoor te zorgen dat de scripts correct functioneren. De vereiste bibliotheken zijn gedefinieerd in het \texttt{requirements.txt}-bestand, zoals weergegeven in codefragment \ref{lst:requirements}. 

\begin{listing} [ht]
    \begin{minted}{text}
        # requirements.txt
        transformers~=4.41.2
        tensorflow~=2.16.2
        torch~=2.3.1
        optimum~=1.20.0
        numpy~=1.26.4
    \end{minted}
    \caption[\texttt{requirements.txt}]{Het \texttt{requirements.txt}-bestand specificeert de versies van bibliotheken die nodig zijn om de scripts te draaien en zorgt voor de juiste compatibiliteit tussen de verschillende Python-pakketten.}
    \label{lst:requirements}
\end{listing}

Het utility-bestand bevat functies die zorgen voor consistent gedrag tussen de verschillende scripts. Deze functies omvatten:

\textbf{Voorbeeldzinnen}: De functie \texttt{get\_sentences} levert vooraf bepaalde zinnen aan elk script, zoals weergegeven in codefragment \ref{lst:getsentences}. Deze zinnen worden gebruikt voor het testen van de modellen en demonstreren correcte en incorrecte grammaticale structuren met betrekking tot de woorden "die" en "dat".

\begin{listing}
    \begin{minted}{python}
        # utils.py
        def get_sentences():
        """Get example sentences."""
        return [
            "Ik heb een vriend die altijd te laat komt.",  # Correct
            "Ik weet die ik het kan.",  # Incorrect
            "Ik weet dat ik het kan.",  # Correct
            "Daarom is het belangrijk, je moet goed opletten.",  # Incorrect
            "Ik ken een man die altijd grapjes maakt.",  # Correct
            "Ze heeft een jurk gekocht die perfect past.",  # Correct
            "Er is een boek dat ik je echt kan aanraden.",  # Correct
            "We bezochten een stad die bekend staat om haar architectuur.",  # Correct
            "Hij las een artikel dat zijn mening veranderde.",  # Correct
            "Ze vertelde over een ervaring die haar leven veranderde.",  # Correct
            "Ik zag een film die mij aan het denken zette.",  # Correct
            "Hij gebruikt een methode die zeer effectief is.",  # Correct
            "Het kind dat in de tuin speelt, is mijn neefje.",  # Correct
            "De vraag die ze stelde, was erg moeilijk.",  # Correct
            "De beslissing dat hij maakte, was erg moeilijk.",  # Incorrect
            "Dit is de laptop die ik wil kopen.",  # Correct
            "Dit is de laptop dat ik wil kopen.",  # Incorrect
            "Het project dat we gestart zijn, verloopt goed.",  # Correct
            "De vrouw die je daar ziet, is mijn lerares.",  # Correct
            "De resultaten dat we behaalden, waren indrukwekkend.",  # Incorrect
            "De resultaten die we behaalden, waren indrukwekkend.",  # Correct
            "De bloemen die je hebt meegenomen, zijn prachtig.",  # Correct
            "De informatie dat hij gaf, was nuttig.",  # Incorrect
            "De informatie die hij gaf, was nuttig.",  # Correct
            "Het idee dat ze voorstelde, was briljant.",  # Correct
            "Het idee die ze voorstelde, was briljant.",  # Incorrect
            "De kat dat op het dak zit, is van ons.",  # Incorrect
            "De kat die op het dak zit, is van ons.",  # Correct
            "De film die we gisteren keken, was spannend.",  # Correct
            "De auto dat ik gisteren zag, was heel mooi.",  # Incorrect
        ]
    \end{minted} 
    \caption[\texttt{get\_sentences()}]{\texttt{get\_sentences()}: Deze functie levert een lijst van voorbeeldzinnen die correct en incorrect gebruik van "die" en "dat" illustreren. Deze zinnen worden gebruikt voor validatie en testing binnen de proof of concept.}
    \label{lst:getsentences}
\end{listing}

\textbf{Preprocessing}: De functie \texttt{preprocess\_input} voert uniforme preprocessing uit voor het genereren van een maskeringstoken binnen een zin, zoals getoond in codefragment \ref{lst:preprocessinput}. Deze functie identificeert de positie van de maskeringstoken en vervangt "die" of "dat" door "<mask>", wat essentieel is voor de verwerking door het model.

\begin{listing}
    \begin{minted}{python}
        # utils.py
        def preprocess_input(sentence):
        """
        Mask 'die' or 'dat' in the input sentence.
        
        Args:
        sentence (str): Input sentence.
        
        Returns:
        tuple: (preprocessed sentence, mask position or None)
        """
        words = sentence.split()
        mask_position = None
        for i, word in enumerate(words):
            if word.lower() in ['die', 'dat']:
                words[i] = '<mask>'
                mask_position = i
                break
        return ' '.join(words), mask_position
    \end{minted}
    \caption[\texttt{preprocess\_input(sentence)}]{\texttt{preprocess\_input(sentence)}: Deze functie voert preprocessing uit op de invoerzinnen door "die" of "dat" te vervangen met "<mask>". Het retourneert zowel de bewerkte zin als de positie van de maskeringstoken.}
    \label{lst:preprocessinput}
\end{listing}

\textbf{Top Voorspellingen}: De functie \texttt{get\_top\_predictions} haalt de top voorspellingen op door \gls{softmax} toe te passen op de \glspl{logit} van het model, zoals weergegeven in codefragment \ref{lst:gettopk}. Dit zorgt voor de selectie van de beste voorspellingen op basis van de maskeringstoken en geeft de top \(k\) voorspelde woorden terug. De \glspl{logit} worden eerst omgezet in \glspl{tensor} voordat \gls{softmax} wordt toegepast, om de waarschijnlijkheden van voorspellingen te verkrijgen.

\begin{listing}
    \begin{minted}{python}
        # utils.py
        def get_top_predictions(logits, vocab, mask_position, top_k=5):
        """
        Get top-k predictions for the masked token.
        
        Args:
        logits (numpy.ndarray or torch.Tensor): Model output logits
        vocab (dict): Vocabulary dictionary
        mask_position (int): Position of the mask token
        top_k (int): Number of top predictions to return
        
        Returns:
        list: Top-k predicted words
        """
        # Convert logits to tensor if it's not already
        if not isinstance(logits, (tf.Tensor, torch.Tensor)):
            logits = tf.convert_to_tensor(logits)
        
        # Get logits for the mask token
        mask_token_logits = logits[0, mask_position + 1]
        
        # Apply softmax to calculate probabilities
        if isinstance(logits, torch.Tensor):
            probs = torch.nn.functional.softmax(mask_token_logits, dim=0)
            top_indices = torch.topk(probs, k=top_k).indices.tolist()
        else:
            probs = tf.nn.softmax(mask_token_logits)
            top_indices = tf.math.top_k(probs, k=top_k).indices.numpy()
        
        # Convert indices to words
        return [list(vocab.keys())[list(vocab.values()).index(idx)] for idx in top_indices]
    \end{minted}
    \caption[\texttt{get\_top\_predictions(logits, vocab, mask\_position, top\_k=5)}]{\texttt{get\_top\_predictions(logits, vocab, mask\_position, top\_k=5)}: Deze functie zet de logit-uitvoer van het model om in \glspl{tensor} en past vervolgens \gls{softmax} toe om de waarschijnlijkheden van voorspellingen te berekenen. Het retourneert de top \(k\) voorspellingen op basis van de maskeringstoken en het vocabulaire.}
    \label{lst:gettopk}
\end{listing}

\subsection{Reproductie met High-Level Pipeline}

In deze fase bespreken we het eerste script waarin \gls{robbert} wordt gebruikt. Zoals eerder vermeld, maken we hierbij gebruik van de high-level \glsxtrshort{api} van \gls{huggingface}, genaamd pipelines. Dit script vormde de basis voor verdere ontwikkeling en testing.

Het script \texttt{main.py}, zoals weergegeven in de codefragmenten \ref{lst:main-part1} en \ref{lst:main-part2}, voert \glsxtrshort{mlm} uit met behulp van de Transformers pipeline op verschillende \gls{robbert}-modellen. Het script maakt gebruik van de \texttt{pipeline} functionaliteit van \gls{huggingface} om modellen eenvoudig te laden en inferentie uit te voeren. Dit vormt een efficiënte manier om gemaskeerde woorden te voorspellen in vooraf bepaalde zinnen.

\begin{listing}
    \begin{minted}[fontsize=\small]{python}
        # main.py
        import logging
        import timeit
        from transformers import pipeline
        from utils import preprocess_input, get_sentences
        
        # Set up logging
        logging.basicConfig(level=logging.INFO)
        
        def set_log_level(log_level):
        """Set the logging level."""
        logging.getLogger().setLevel(log_level)
        
        def run_inference(sentence, pipe):
        """Process sentence with mask, predict masked word."""
        if "<mask>" not in sentence:
        logging.debug("No mask token found in the sentence.")
        return
        return pipe(sentence)[0]
        
        def main(model_path):
        """Main function to process sentences with the specified model."""
        logging.info("========================================")
        
        # Load the model
        pipe = pipeline("fill-mask", model=model_path)
        logging.info(f"Model loaded: {model_path}")
        
        # Example sentences
        sentences = get_sentences()
        
        start_total_time = timeit.default_timer()
        for sentence in sentences:
        masked_sentence, _ = preprocess_input(sentence)
        start_time = timeit.default_timer()
        output = run_inference(masked_sentence, pipe)
        elapsed_time = timeit.default_timer() - start_time
        logging.info(f"Original sentence: {sentence}")
        logging.info(f"Processed sentence: {masked_sentence}")
        logging.info(f"Prediction: {output}")
        logging.info(f"Elapsed time: {elapsed_time:.3f} seconds\n")
        logging.info("")
        
        total_time = timeit.default_timer() - start_total_time
        logging.info(f"Total elapsed time: {total_time:.3f} seconds")
        logging.info("Average time per sentence: {:.3f} seconds".format(total_time / len(sentences)))
    \end{minted}
    \caption[\texttt{main.py} (deel 1)]{Het eerste deel van \texttt{main.py} definieert de functies die nodig zijn voor het uitvoeren van \glsxtrshort{mlm}. Het script laadt modellen en verwerkt zinnen met behulp van de \texttt{pipeline} functionaliteit van \gls{huggingface}. De belangrijkste functies zijn \texttt{run\_inference}, die voorspellingen doet voor gemaskeerde woorden, en \texttt{main}, die de zinnen ophaalt, modellen laadt, en de totale verwerkingstijd registreert.}
    \label{lst:main-part1}
\end{listing}

\begin{listing}
    \begin{minted}[fontsize=\small]{python}
        # main.py
        if __name__ == '__main__':
        """
        This script performs masked language modeling using the Transformers pipeline
        on various Dutch BERT models. It provides a simple way to load and run inference
        using the high-level pipeline API.
        
        Key features:
        1. Uses Transformers pipeline for easy model loading and inference
        2. Supports multiple Dutch BERT models
        3. Preprocesses input sentences to include a mask token
        4. Runs inference to predict masked words
        5. Provides logging of original sentences, processed sentences, and predictions
        
        Usage:
        Run the script to process a set of predefined sentences with different Dutch
        BERT models using the Transformers pipeline.
        
        Models used:
        - DTAI-KULeuven/robbertje-1-gb-non-shuffled5
        - pdelobelle/robbert-v2-dutch-base
        - DTAI-KULeuven/robbert-2022-dutch-base
        - DTAI-KULeuven/robbert-2023-dutch-base
        - DTAI-KULeuven/robbert-2023-dutch-large
        """
        
        # Uncomment to set debug logging
        # set_log_level(logging.DEBUG)
        
        # Process sentences with different models
        main("DTAI-KULeuven/robbertje-1-gb-non-shuffled")
        main("pdelobelle/robbert-v2-dutch-base")
        main("DTAI-KULeuven/robbert-2022-dutch-base")
        main("DTAI-KULeuven/robbert-2023-dutch-base")
        main("DTAI-KULeuven/robbert-2023-dutch-large")
    \end{minted}
    \caption[\texttt{main.py} (deel 2)]{Het tweede deel van \texttt{main.py} voert de hoofdtoepassing uit, waarbij gebruik wordt gemaakt van een lijst met vooraf gedefinieerde modellen. Dit gedeelte initialiseert het script door verschillende \gls{robbert}-modellen te testen en de inferentie uit te voeren op de vooraf bepaalde zinnen. Het documenteert de kernfuncties en gebruiksinstructies in de \texttt{docstring} voor verduidelijking en instructies.}
    \label{lst:main-part2}
\end{listing}

Het script begint met het instellen van logging via de \texttt{logging} module om inzicht te geven in de verwerkingsstappen en resultaten. De functie \texttt{set\_log\_level} maakt het mogelijk om het logniveau aan te passen, wat nuttig is voor het controleren van gedetailleerde uitvoer tijdens het debuggen.

De kern van het script is de \texttt{main} functie, die verantwoordelijk is voor het verwerken van zinnen met een gespecificeerd model. Binnen deze functie wordt eerst het model geladen met behulp van de \texttt{pipeline} \glsxtrshort{api} van \gls{huggingface}, wat het eenvoudig maakt om verschillende modellen te gebruiken door enkel het pad naar het model te specificeren.

Vervolgens haalt de functie \texttt{get\_sentences} een lijst op van voorbeeldzinnen die zullen worden gebruikt voor inferentie. Elke zin wordt voorbewerkt met de \texttt{preprocess\_input} functie, die het masker token in de zin plaatst en de positie ervan identificeert. Dit is cruciaal voor \gls{robbert} om de juiste voorspellingen te kunnen maken.

De functie \texttt{run\_inference} neemt de gemaskeerde zin en het geladen model om het gemaskeerde woord te voorspellen. Het controleert of er een \texttt{<mask>} token aanwezig is; zo niet, dan wordt een debugbericht gelogd en wordt de functie beëindigd zonder inferentie. Voor zinnen met een \texttt{<mask>} token retourneert het de topvoorspelling.

Het script meet de tijd die nodig is om elke zin te verwerken, evenals de totale verwerkingstijd voor alle zinnen. Dit gebeurt met de \texttt{timeit} module, die helpt bij het vastleggen van prestatiemetrics die nuttig zijn voor analyse.

Tot slot biedt het script flexibiliteit door meerdere modellen te ondersteunen. Dit maakt het mogelijk om de prestaties van verschillende modellen te vergelijken en te testen binnen dezelfde structuur. De uitvoer wordt gelogd, inclusief de originele en gemaskeerde zinnen, de voorspelde woorden en de benodigde tijd voor elke inferentie.

Dit script is een essentieel onderdeel van het proof of concept, aangezien het de mogelijkheden van \gls{robbert}-modellen demonstreert met behulp van de high-level pipeline \glsxtrshort{api} van \gls{huggingface}. Het is de basis voor verdere ontwikkeling en evaluatie van de komende integraties. Houd er rekening mee dat dit script een internetverbinding vereist om de \glsxtrshort{api} te kunnen gebruiken.

\subsection{Transformer Bibliotheek: Low-Level API}
In dit script maken we gebruik van de Transformer-bibliotheek van \gls{huggingface} op een lower-level. Dit script verschilt enigszins van het eerste script, dat de high-level \glsxtrshort{api} via pipelines benutte. Door over te stappen naar een lower-level benadering, krijgen we beter inzicht in de interne werking van de \gls{huggingface} pipeline, wat de overgang naar TensorFlow Lite (\glsxtrshort{tflite}) vergemakkelijkt. 



\subsection{Conversie naar TFLite}
Het doel van deze stap is om de RobBERT-modellen te converteren naar het \glsxtrshort{tflite}-formaat met behulp van \gls{huggingface} Optimum. Deze conversie maakt het mogelijk om de modellen efficiënt te implementeren op mobiele apparaten. Bovendien biedt deze stap de mogelijkheid om gekwantiseerde modellen als output te genereren, wat de prestaties verder kan verbeteren.



\subsection{Implementatie van TFLite Model in Python}
Deze sectie beschrijft de implementatie van de geconverteerde \glsxtrshort{tflite}-modellen in een Python-script dat de \glsxtrshort{tflite} interpreter \glsxtrshort{api} gebruikt. Het script voert inferenties uit met deze modellen en maakt gebruik van een aangepaste tokenizer voor de verwerking van invoertekst.




\section{Android Implementatie}

Alle code en functionaliteit die in deze sectie wordt besproken, is beschikbaar in de volgende GitHub-repository: \url{https://github.com/Chmpy/RobBERT_local_Android_example}. Deze repository de Android-applicatie in zijn geheel voor het integreren van \gls{robbert}

\subsection{Compatibiliteit en Instellingen}

In deze sectie bespreken we de specificaties van de Android-applicatie om compatibiliteit en consistentie met de eerder genoemde scripts en modellen te waarborgen. Het utility-bestand is van cruciaal belang voor het handhaven van uniform gedrag van de \glsxtrshort{tflite}-modellen binnen de applicatie. Daarnaast bevat het manifestbestand de nodige instellingen om de applicatie correct te laten functioneren op mobiele apparaten. Deze elementen zorgen ervoor dat de applicatie consistent presteert in overeenstemming met de Python-implementaties.



\subsection{Ontwikkeling van de Android-applicatie}

Dit gedeelte beschrijft de ontwikkeling van de Android-applicatie, met de nadruk op de integratie van het \glsxtrshort{tflite}-model. De eerder ontwikkelde Python-code wordt omgezet voor gebruik op Android, met behulp van de \glsxtrlong{tflite} Android-bibliotheek om het model te integreren. De \texttt{MainActivity} fungeert als het centrale onderdeel van de applicatie, waar de inferentie met het \glsxtrshort{tflite}-model wordt uitgevoerd. De inferentieprocedure is consistent met die in de Python-implementatie, met aanvullende functionaliteiten voor Android-specifieke taken.



\subsection{Integratie met Process Text Intent}

Als laatste binnen dit gedeelte behandelen we de implementatie van de \gls{processtextintent} binnen de Android-applicatie. Deze functionaliteit maakt het mogelijk om inferenties uit te voeren op tekst, zowel direct binnen de applicatie als via externe intents. De inferentieprocedure is gelijk aan die in de \texttt{MainActivity}, maar de gegevensverwerking verschilt enigszins vanwege het gebruik van Android-specifieke methoden. Deze aanpak verzekert dat de applicatie flexibel is en naadloos samenwerkt met andere apps op het Android-platform.




\section{Evaluatie en Vergelijking van Inferentie}

Nu we alle code en implementaties hebben besproken, kunnen we overgaan tot de evaluatie en vergelijking van de prestaties van alle scripts en de Android-applicatie. In deze sectie richten we ons op de verschillen of overeenkomsten in hun prestaties en geven we inzicht in waarom deze prestatieverschillen mogelijk bestaan. We zullen enkele belangrijke bevindingen en relevante metingen delen om inzicht te geven in de effectiviteit van elk script en de applicatie.


